\bookchapter{บทที่ 36: มือถือกับสมาร์ทโฮม - ทำให้บ้านฉลาดขึ้น}

\chapterintro{มือถือยังทำอีกเรื่องที่สำคัญขึ้นทุกที คือเป็น “รีโมต” ของทั้งบ้าน ตั้งแต่ดวงไฟดวงเดียว ปลั๊กไฟสักตัว ไปจนถึงกุญแจประตู สมาร์ทโฮมกำลังเปลี่ยนจาก “ของเล่นคนสายเทค” มาเป็น “อุปกรณ์มาตรฐานในบ้าน” บทนี้จะช่วยคุณไล่เรียงเรื่องโปรโตคอล ระบบนิเวศ ฮับกลาง และชุดเริ่มต้น แล้วปิดท้ายด้วยเช็กลิสต์ความปลอดภัย เพื่อให้คุณไม่ถูกความ “อัจฉริยะ” ย้อนกลับมาทำร้ายเอง}

\booksection{โปรโตคอลสมาร์ทโฮม: การเลือกมีความสำคัญ}

ก่อนซื้ออุปกรณ์สมาร์ท ให้ดูก่อนว่ามันรองรับโปรโตคอลอะไร เพราะสิ่งนี้ตัดสินว่ามือถือและระบบนิเวศของคุณจะ “สั่งงาน” มันได้หรือเปล่า ตอนนี้มีโปรโตคอลหลักอยู่สี่ตัว นิสัยของแต่ละตัวเป็นดังนี้:

\begin{itemize}
  \item โปรโตคอล Matter: เปิดตัวปี 2022 โดย CSA (Connectivity Standards Alliance) โดยมียักษ์ใหญ่อย่าง Apple, Google, Amazon, Samsung ร่วมมือกันอยู่เบื้องหลัง เป้าหมายมีอย่างเดียว คือทำให้อุปกรณ์สมาร์ททุกแบรนด์พูดภาษาเดียวกัน ช่วงปี 2023-2026 มันกำลังกลายเป็นมาตรฐานเดียวกันอย่างรวดเร็ว อุปกรณ์ที่รองรับจะเชื่อมเข้า Apple HomeKit, Google Home, Amazon Alexa และ Mi Home ได้พร้อมกัน ไม่ต้องโดนผูกขาดอยู่ในระบบนิเวศเดียว
  \item โปรโตคอล Thread: ท่อส่งเครือข่ายระดับล่างของ Matter กินไฟน้อย จัดเครือข่ายเองได้ (ถ้าโหนดไหนไฟดับ สัญญาณจะอ้อมเองอัตโนมัติ) ตอบสนองไว ไม่สะดุด ซื้ออุปกรณ์ที่มี Thread Border Router ในตัว (เช่น HomePod mini, Apple TV 4K, เกตเวย์ฮับกลางของ Xiaomi) ประสบการณ์จะนิ่งกว่าชัดเจน
  \item Wi-Fi: แบนด์วิดท์กว้าง เชื่อมตรงเข้าเราเตอร์ได้ แต่กินไฟเยอะ ถ้าอุปกรณ์เยอะจะแย่งช่องสัญญาณง่าย เหมาะกับกล้องวงจรปิด หุ่นยนต์ดูดฝุ่น ฯลฯ ที่ต้องส่งข้อมูลปริมาณใหญ่
  \item Bluetooth (BLE): ตั้งค่าเชื่อมเน็ตง่าย ประหยัดไฟ แต่ระยะสั้น และจัดเครือข่ายเองไม่ได้ เหมาะกับอุปกรณ์ข้อมูลเล็ก ๆ อย่างกุญแจดิจิทัลหรือเครื่องวัดอุณหภูมิ-ความชื้น
  \item Zigbee: โปรโตคอลกินไฟน้อยจัดเครือข่ายเองได้รุ่นเก๋า ต้องมีเกตเวย์เฉพาะ อุปกรณ์ถูกและเสถียร แต่ความเข้ากันได้ข้ามแบรนด์แย่ และกำลังถูก Matter+Thread แทนที่ทีละนิด
\end{itemize}

\begin{tcbinfo}{หลักการแบบหนึ่งประโยค}
ซื้ออุปกรณ์สมาร์ทปี 2026 ให้ดูก่อนว่ามีตราสัญลักษณ์ “Matter over Thread” ไหม นี่คือทางเลือกที่ผ่านการพิสูจน์เวลาได้ดีที่สุดตอนนี้ อุปกรณ์เก่าที่รองรับแค่ Wi-Fi หรือ Bluetooth ก็ไม่ต้องรีบเลิกใช้ ใช้ต่อไปก่อน ค่อยเปลี่ยนมา Matter ตอนซื้อใหม่

\end{tcbinfo}

\booksection{เปรียบเทียบสามระบบนิเวศใหญ่}

โปรโตคอลคือ “ภาษา” ส่วนระบบนิเวศคือ “วง” แม้ทุกอย่างจะรองรับ Matter ก็ตาม คุณยังต้องเลือกหนึ่งระบบนิเวศหลักไว้ใช้ทำระบบอัตโนมัติและควบคุมด้วยเสียงในชีวิตประจำวัน สามระบบนิเวศหลักเปรียบเทียบกันดังนี้:

\begin{center}
\begin{tabularx}{\linewidth}{|X|X|X|X|}
\hline
\textbf{ระบบนิเวศ} & \textbf{จุดแข็ง} & \textbf{จุดอ่อน} & \textbf{เหมาะกับใคร} \tabularnewline
\hline
Apple HomeKit & ความเป็นส่วนตัวแน่นที่สุด (ประมวลผลในเครื่องเป็นหลัก) ทำงานร่วมกับ iPhone/Siri/Shortcuts ลื่นไหล & อุปกรณ์น้อยและราคาสูง ต้องมี HomePod หรือ Apple TV เป็นฮับกลาง & ครัวเรือนที่ใช้ของ Apple ทั้งบ้าน ผู้ใช้ที่ให้ความสำคัญกับความเป็นส่วนตัวและประสบการณ์ \tabularnewline
Mi Home (Xiaomi) & อุปกรณ์เยอะที่สุด คุ้มค่าที่สุด และผู้ช่วยเสียง Xiao AI ใช้งานได้ดี & ระบบนิเวศค่อนข้างปิด ข้อมูลส่วนใหญ่เก็บบนคลาวด์ โฆษณาก็เยอะหน่อย & ผู้เริ่มต้นที่งบจำกัด อยากจัดเต็มทั้งบ้านแบบจบในที่เดียว \tabularnewline
Huawei AI Life (HarmonyOS) & ส่งต่องานระหว่างมือถือ/แท็บเล็ต/นาฬิกา HarmonyOS ได้ไร้รอยต่อ จำนวนอุปกรณ์โตเร็ว & อุปกรณ์สากล/ต่างประเทศรองรับน้อย ระบบนิเวศยังไล่ตามอยู่ & ผู้ใช้มือถือ Huawei คนที่ให้ความสำคัญกับ "Super Device" ที่ทำงานร่วมกันหลายอุปกรณ์ \tabularnewline
\hline
\end{tabularx}
\end{center}

\begin{tcbinfo}{คำแนะนำตัดสินใจเลือกระบบนิเวศ}
เลือกระบบนิเวศจากมือถือที่ใช้: ผู้ใช้ iPhone เอา HomeKit ก่อน ผู้ใช้ Huawei เอา AI Life ก่อน ผู้ใช้ Android คนอื่น Mi Home สบายใจที่สุด และด้วยพลังของ Matter อุปกรณ์คนละระบบนิเวศก็ควบคุมได้ด้วยมือถือเครื่องเดียว ไม่ต้อง “เลือกข้าง” ตายตัวตั้งแต่แรก

\end{tcbinfo}

\booksection{ใช้มือถือเป็นฮับกลางของบ้าน}

ในกรณีใช้งานระดับเริ่มต้นส่วนใหญ่ มือถือก็คือ “ฮับกลางของบ้าน” ของคุณแล้ว ควบคุมไฟ แอร์ ม่าน และระบบรักษาความปลอดภัยผ่านแอป ยังสร้างฉากอัตโนมัติได้ ให้บ้าน “ขยับเองได้”:

\begin{itemize}
  \item ควบคุมอุปกรณ์: เปิดแอป HomeKit/Mi Home/AI Life สามารถเปิด-ปิดไฟ ปรับเฉดสีแสง ตั้งอุณหภูมิแอร์ ดึงม่าน และดูภาพจากกล้อง ได้ทั้งจากระยะไกลและในเครือข่ายเดียวกัน
  \item ฉากอัตโนมัติ: จับ "เงื่อนไข การทำงาน" มาประกอบกัน ฉากหลักที่ใช้บ่อยมีสามแบบ: โหมดกลับบ้าน: ตำแหน่งถึงบ้าน+พระอาทิตย์ตก เปิดไฟหน้าประตู+เปิดแอร์+ดึงม่าน+ล็อกประตู
  \item โหมดออกจากบ้าน: คนสุดท้ายออกจากบ้าน (ตำแหน่งออกนอกบ้านหรือประตูถูกล็อก) ปิดไฟ+ปิดแอร์+ปิดปลั๊ก+เปิดกล้อง+ตั้งระบบรักษาความปลอดภัย
  \item โหมดนอน: ทุกคืน 23:00 ปิดไฟหลักทั้งบ้าน+ดึงม่าน+แอร์ปรับ 26 องศา+เปิดไฟนอน
\end{itemize}

NFC แตะทีเดียว: แปะแท็ก NFC ที่ประตู แล้วเอามือถือแตะเบา ๆ ก็เรียกฉากที่ตั้งไว้ได้ (แตะครั้งแรกเปิดไฟหน้าประตู แตะอีกครั้งสลับเข้าโหมดนอน) iPhone ใช้ Shortcuts ทำอัตโนมัติ ฝั่ง Android ใช้แอป Wallet หรือแอปผู้ผลิตเขียนแท็กได้ทั้งนั้น

\begin{tcbpractice}{มาลองดู: สร้าง "โหมดกลับบ้าน" หนึ่งอัน}
ในแอปสมาร์ทโฮมของมือถือ เข้าเมนู “อัตโนมัติ” แล้วสร้างรายการใหม่: เงื่อนไขทริกเกอร์เลือก “มาถึงสถานที่” (บ้าน) แล้วเพิ่มการทำงานทีละอย่าง “เปิดไฟหน้าประตู” “ตั้งแอร์ 26 องศา” “ปิดม่าน” บันทึกแล้วออกไปเดินเล่นรอบหนึ่ง กลับมาถึงบ้านดูซิว่าไฟเปิดเองไหม ความภูมิใจครั้งแรกที่สำเร็จ ฟินกว่ากดสั่งซื้ออุปกรณ์ใหม่อีก

\end{tcbpractice}

\booksection{ชุดเริ่มต้นที่แนะนำ}

ไม่ต้องจัดให้ครบในคราวเดียว เริ่มจากสี่ชิ้นนี้ก็ครอบคลุมความต้องการประจำวันได้แปดสิบเปอร์เซ็นต์แล้ว:

\begin{itemize}
  \item ลำโพงอัจฉริยะ: Xiao AI (ระบบนิเวศ Mi Home คุ้มค่า) / HomePod mini (ระบบนิเวศ Apple ทำหน้าที่ Thread hub ด้วย) / Alexa (ระบบนิเวศสากลกว้าง) ระวัง: ลำโพงตัวนี้จะตัดสินระบบนิเวศเสียงของคุณ เลือกระบบนิเวศก่อนแล้วค่อยซื้อลำโพง
  \item หลอดไฟอัจฉริยะ: Yeelight (ระบบ Mi Home/Matter ปรับเฉดสีแสงและความสว่างได้ ราคาถูก) / Philips Hue (สีสวย นิ่ง แต่ต้องมี Hue Bridge และราคาสูงหน่อย) ระวัง: ก่อนซื้อเช็กขนาดขั้วหลอด (E27/B22) และว่าต้องปรับเฉดสีไหมด้วย
  \item ปลั๊กอัจฉริยะ: Xiaomi Smart Plug (ถูก มีสถิติการใช้ไฟ) / TP-Link Kasa (รองรับ Matter สำหรับตลาดสากลดี) ระวัง: ถ้าใช้กับเครื่องใช้ไฟกำลังสูง (เครื่องทำน้ำอุ่น เครื่องทำความร้อน) ต้องดูกำลังไฟที่รับได้ให้ดี อย่าให้โอเวอร์โหลด
  \item กุญแจดิจิทัล: Loock/Xiaomi (คุ้มค่า เชื่อมกับ Mi Home ได้) / Kaadas (ผู้ผลิตล็อกรายใหญ่สายดั้งเดิม โครงสร้างกลไกแน่น) ระวัง: ยืนยันว่ารองรับ Matter หรือระบบนิเวศที่คุณใช้ไหม และสังเกตอายุแบตเตอรี่กับตำแหน่งรูกุญแจฉุกเฉินด้วย
\end{itemize}

\begin{tcbinfo}{ลำดับการซื้อที่แนะนำ}
ซื้อลำโพงอัจฉริยะก่อน (ตัดสินระบบนิเวศ) ต่อด้วยหลอดไฟและปลั๊กอัจฉริยะ (เริ่มง่าย เห็นผลทันที) ท้ายสุดคือกุญแจดิจิทัลและกล้อง (เกี่ยวกับความปลอดภัย เลือกอย่างรอบคอบ) เพิ่มทีละประเภท พอคุ้นแล้วค่อยขยาย กันพลาดเรื่อง “ซื้อมาแล้วใช้ไม่เป็น”

\end{tcbinfo}

\booksection{เรื่องความปลอดภัยที่ต้องระวัง}

บ้านไหนอุปกรณ์อัจฉริยะเยอะ พื้นที่โจมตียิ่งกว้าง สี่ข้อต่อไปนี้คือเส้นตาย ลองเช็กทีละข้อ:

\begin{enumerate}
  \item เปลี่ยนรหัสผ่านเริ่มต้น: รหัสผ่านจากโรงงานของกล้อง เราเตอร์ และกุญแจดิจิทัลหาได้ทั้งบนอินเทอร์เน็ต พอได้ของมาถึงมือ สิ่งแรกที่ต้องทำคือเปลี่ยนเป็นรหัสแข็งแรงแบบ "ยาวและไม่ซ้ำกัน"
  \item อัปเดตเฟิร์มแวร์ให้ทัน: เฟิร์มแวร์ที่ผู้ผลิตปล่อยมักมาพร้อมแก้ช่องโหว่ เปิด "อัปเดตอัตโนมัติ" ในแอป หรืออย่างน้อยก็เข้าไปเช็กเองเดือนละครั้ง
  \item แยกเครือข่าย: พาอุปกรณ์ IoT ไปอยู่ Wi-Fi หนึ่ง (หรือเครือข่ายแขก) แยกจากเครือข่ายหลักของมือถือและคอม ถ้าอุปกรณ์ตัวไหนถูกเจาะ ก็จะลามไปขู่กระแทกข้อมูลสำคัญของคุณได้ยากขึ้น
  \item ปกป้องความเป็นส่วนตัว: เปิด "โซนส่วนตัว" ให้กล้องปิดบัง (เช่น ห้องนอน ห้องน้ำไม่ต้องถ่าย) ปิดการอัดวิดีโอเก็บบนคลาวด์ 24 ชั่วโมง และอย่าเปิดภาพกล้องตอนไม่อยู่บ้านให้แอปบุคคลที่สามใช้ยาว ๆ
\end{enumerate}

\textbf{หลุมพรางพวกนี้อย่าเหยียบ}\par

\begin{itemize}
  \item อย่าใช้ฟังก์ชัน "ปลดล็อกระยะไกล" ของกุญแจให้คนแปลกหน้า เลือกใช้รหัสชั่วคราวแทน ใช้เสร็จก็หมดอายุทันที
  \item อย่าหันกล้องเฉียบตรงเข้าเตียงหรือออกนอกหน้าต่าง ของที่รั่วไหลไปคือความเป็นส่วนตัวของคุณเอง
  \item อย่าเห็นแก่ถูกซื้ออุปกรณ์ไม่มีแบรนด์แบบ "white label" พวกนี้ส่วนใหญ่ไม่อัปเดตเฟิร์มแวร์ เท่ากับติดประตูหลังลับไว้ที่บ้าน
\end{itemize}

\begin{tcbsummary}{สรุปบทนี้}
\begin{itemize}
  \item Matter+Thread คือทิศทางรวมศูนย์ของปี 2023-2026 ซื้ออุปกรณ์ใหม่ให้ดูตราสัญลักษณ์นี้ก่อน
  \item สามระบบนิเวศต่างมีจุดเด่นของตัวเอง: Apple เน้นความเป็นส่วนตัวและประสบการณ์, Mi Home เน้นความคุ้มค่า, Huawei เน้นการทำงานร่วมกันหลายอุปกรณ์
  \item มือถือเป็นฮับกลางได้ ตั้งสามฉากใหญ่ "กลับบ้าน/ออกจากบ้าน/นอน" ให้เป็น บ้านก็ขยับเองได้
  \item ชุดเริ่มต้นสี่ชิ้น: ลำโพงอัจฉริยะ หลอดไฟ ปลั๊ก กุญแจดิจิทัล เพิ่มทีละอย่างตามลำดับนี้
  \item เปลี่ยนรหัสเริ่มต้น อัปเดตเฟิร์มแวร์ แยกเครือข่าย บังพื้นที่ส่วนตัว - สี่เส้นตายด้านความปลอดภัยห้ามละเลย
\end{itemize}

\end{tcbsummary}
